\chapter{Dec.~1 --- Representations of \texorpdfstring{$\GL_n(\C)$}{GL(n, C)}}

\section{More on Irreducible Representations}
\begin{prop}
  Let $\g$ be a semisimple Lie algebra
  and $\lambda = \sum_{i = 1}^r m_i \omega_i$
  be a dominant integral weight.
  Let $T_\lambda = \bigotimes_i L_{\omega_i}^{\otimes m_i}$,
  $v = \bigotimes_i v_i^{\otimes m_i}$,
  and $V$ the subrepresentation
  of $T_\lambda$ generated by $v$.
  Then $V \cong L_\lambda$.
\end{prop}

\begin{proof}
  Let $C = \sum_j x_j^2 + \sum_{\alpha \in R_+} (e_{-\alpha} e_\alpha + e_\alpha e_{-\alpha})$
  be the Casimir element, where
  $x_j$ is an orthonormal basis
  of the Cartan subalgebra
  $\mathfrak{h} \subset \g$. Since
  $[e_\alpha, e_{-\alpha}] = h_\alpha$,
  we can write
  \[
    C = \sum_j x_j^2 + 2\sum_{\alpha \in R_+} e_{-\alpha} e_\alpha
    + \sum_{\alpha \in R_+} h_\alpha.
  \]
  Now note the following:
  \begin{quote}
    \vspace{-2em}
    \begin{lemma}
      $C|_V = (\lambda, \lambda + 2\rho) = |\lambda + \rho|^2 - |\rho|^2$,
      where $\rho = \frac{1}{2} \sum_{\alpha \in R_+} \alpha$ and
      $V$ is the highest weight representation
      with highest weight $\lambda$.
    \end{lemma}
  \end{quote}
  Write
  $V = L_\lambda \oplus \bigoplus_{\mu \in (\lambda - Q_+) \cap P_+} N_{\lambda, \mu} L_\mu$,
  where $N_{\lambda, \mu} \in \Z$. We
  know
  \[
    C|_{L_\mu} = (\mu, \mu + 2\rho)
    \quad \text{and} \quad
    C|_V = (\lambda, \lambda + 2\rho).
  \]
  One can check as an exercise that
  for every $\mu \in (\lambda - Q_+) \cap P_+$
  such that $\mu \ne \lambda$, we have
  $(\mu, \mu + 2\rho) < (\lambda, \lambda + 2\rho)$.
  Thus we must have $N_{\lambda, \mu} = 0$.
\end{proof}

\section{Representations of \texorpdfstring{$\SL_n(\C)$}{SL(n, C)}}

\begin{remark}
  We now consider representations of
  $\SL_n(\C)$, which are the
  same as the representations of
  $\mathfrak{sl}_n(\C)$.
  Let $\lambda = \sum_{i = 1}^{n - 1} m_i \omega_i$
  (the $\omega_i$ are the fundamental
  weights) and $L_\lambda$ the
  corresponding representation.
  Then a Cartan subalgebra $\mathfrak{h} \subseteq \mathfrak{sl}_n(\C)$
  is the space of $\C_0^n$ vectors
  (i.e. those with zero sum of
  coordinates), and
  $\mathfrak{h}^* \subseteq \C$
  are vectors $(x_1, \dots, x_n)$ which
  act by (modulo) shifting by the same
  number. Note that
  $\mathfrak{h}^* = \C^n / \C_{\mathrm{diag}}$.
  Then
  $\alpha_i^\vee = e_i - e_{i + 1}$, and
  we have
  $(\omega, e_j - e_{j + 1}) = \delta_{i, j}$
  Then
  \[
    \omega_i = (\underbrace{1, \dots, 1}_{i \text{ times}}, 0, \dots, 0)
  \]
  The dominant integral weight is
  given by $\lambda = (m_1 + \dots + m_{n - 1}, m_2 + \dots + m_{n - 1}, \dots, m_{n - 1}, 0)$.
  Note that we can write
  $\lambda = (\lambda_1, \dots, \lambda_{n - 1}, 0)$
  where $\lambda_1 \ge \dots \ge 0$.

  Now consider the vector representation
  of $\mathfrak{sl}(\C)$. Let
  $e = E_{i, i + 1}$,
  $v_1, \dots, v_n$ be the standard basis,
  so that
  $v_{1} 
  \wedge \cdots \wedge v_{k}$ is the
  highest weight vector.
\end{remark}

\begin{remark}
  We have the following:
  \begin{enumerate}
    \item $\wedge^n V \cong \C$ is
      the trivial representation;
    \item $V^* \cong \wedge^{n - 1} V$ and
      $V \otimes \wedge^{n - 1} \to \wedge^n  V \cong \C$.
      Similarly,
      $\wedge^m V^* \cong \wedge^{n - m} V$.
  \end{enumerate}
\end{remark}

\section{Representations of \texorpdfstring{$\GL_n(\C)$}{GL(n, C)}}

\begin{remark}
  To go from $\SL_n(\C)$ to
  $\GL_n(\C)$, we can write
  \[
    \GL_n(\C)
    = (\C^\times \times \SL_n(\C)) / \mu_n,
  \]
  where $\mu_n$ is the group of
  $n$th roots of unity.
  We have $z \mapsto (z^{-1}, z 1_n)$
  and $(\lambda, B) \sim (z^{-1} \lambda, z B)$.

  Thus irreducible holomorphic representations
  of $\GL_n(\C)$ is a subset of the
  irreducible holomorphic representations
  of $\C^\times \times \SL_n(\C)$.
\end{remark}

\begin{example}
  For $n = 1$, we just have
  $\C^\times$. Then $e^{2\pi i h} = 1$,
  where $h$ is a generator of the Lie
  algebra, and $h \mapsto H$
  gives $e^{2\pi i H} = 1$.
  Now consider the character $\chi_N : z \mapsto z^N$ for
  $N \in \Z$.

  If we look at $\C^\times \times \SL_n(\C)$,
  let $L_{\lambda, N} = \chi_N \otimes L_\lambda$.
  For $\GL_n$, consider $L_\lambda $ and
  take
  \[
    N = n m_n + \sum_{i = 1}^{n - 1} \lambda_i.
  \]
  Then we get
  \[
    (\lambda_1, \dots, \lambda_n)
    = (m_1 + \dots + m_n, m_2 + \dots + m_n, \dots, m_{n - 1} + m_n, m_n),
  \]
  where $m_n$ is not necessarily positive.
  This corresponds to
  $\wedge^n V$, which is a nontrivial
  $1$-dimensional representation.
  Here $\omega_n = (1, 1, \dots, 1)$.
  Let $L_\lambda \subseteq \bigotimes_{i = 1}^n (\wedge^i V)^{\otimes m_i}$,
  where we interpret
  $\chi^{\otimes k} = (\chi^*)^{\otimes (-k)}$
  for $m_n < 0$.
  When we do have $m_n \ge 0$, the
  representation is called
  \emph{polynomial}.
\end{example}

\begin{remark}
  Every polynomial representation occurs
  inside $V^{\otimes N}$, so
  we get a natural action by $S_N$.
\end{remark}

\section{Schur-Weyl Duality}

\begin{remark}
  Consider the highest weight of
  a polynomial representation of
  $|\lambda| = \sum_i \lambda_i$, where
  the $\lambda_i$ are the eigenvalues
  of $1_n \in \gl_n(\C)$.
  Then $|\lambda| = N$, so we see that
  \[
    \lambda_1, \dots, \lambda_n
  \]
  is a partition with $\le n$ parts.
  Now recall Young diagrams. We can write
  \[
    V^{\otimes N}
    = \bigoplus_{\lambda : |\lambda| = N} L_\lambda \otimes \pi_\lambda,
  \]
  where the $\pi_\lambda = \Hom_{\GL_n(\C)}(L_\lambda, V^{\otimes N})$ are multiplicity
  spaces. Note that
  $L_\lambda = 0$ if the partition
  has more than $n$ parts. The natural
  action of $S_N$ on $V^{\otimes N}$
  commutes with the $\GL_n(\C)$ action,
  so it acts on the multiplicity spaces
  $\pi_\lambda$. Let
  $A$ be the image of
  $U(\gl_n(\C))$ in $\End_\C(V^{\otimes N})$
  and $B$ the image of
  $\C S_N$.
\end{remark}

\begin{theorem}[Schur-Weyl duality]\label{thm:schur-weyl}
  Let $A$ and $B$ as above.
  \begin{enumerate}
    \item The centralizer of $A$ is
      $B$ and vice versa.
    \item If $\lambda$ has at most
      $n$ parts, then the representation
      $\pi_\lambda$ of $B$ (and therefore
      of $S_N$) is irreducible, and
      such representations are pairwise
      isomorphic.
    \item If $\dim V \ge N$, then
      the $\pi_\lambda$ exhaust
      all irreducible representations of
      $S_N$.
  \end{enumerate}
\end{theorem}

\begin{lemma}\label{lem:symmetric-span}
  If $U$ is a $\C$-vector space, then
  $S^N U$ is spanned by elements
  $x \otimes \cdots \otimes x$
  for $x \in U$.
\end{lemma}

\begin{proof}
  For simplicity we assume $U$ is
  finite-dimensional. Then the
  span of such vectors generates a
  nonzero subrepresentation of the
  irreducible $\GL(U)$
  representation $S^N U$, so the
  statement follows.
\end{proof}

\begin{lemma}
  For any associative algebra
  $R$ over $\C$, the algebra
  $S^N R$ is generated by elements
  \[
    \Delta_N(x) =
    (x \otimes 1 \otimes \cdots \otimes 1)
    + (1 \otimes x \otimes 1 \otimes \cdots \otimes 1)
    + \cdots
    + (1 \otimes \cdots \otimes 1 \otimes x).
  \]
\end{lemma}

\begin{proof}
  We can write
  $z_1 \cdots z_N = P_N(p_1, \dots, p_N)$
  for some polynomial $P_N$ and
  $p_k = \sum_i z_i^k$ for $k = 1, \dots, N$.
  Thus we have
  \[
    x \otimes \cdots \otimes x
    = P_N(\Delta_1(x), \dots, \Delta_N(x)),
  \]
  so the statement follows from
  Lemma~\ref{lem:symmetric-span}.
\end{proof}

\begin{lemma}[Double centralizer theorem]\label{lem:double-centralizer}
  Let $V$ be a finite-dimensional
  vector space and $A, B \subseteq \End(V)$ 
  be subalgebras. Assume $B$ is
  isomorphic to a direct sum of matrix
  algebras and
  $A$ is a centralizer of $B$.
  Then $A$ is also isomorphic to a direct
  sum of matrix algebras, and
  \[
    V = \bigoplus_{i = 1}^n W_i \otimes U_i,
  \]
  where the $W_i$ run through all irreducible
  $A$-modules, the $U_i$ run through
  all irreducible $B$-modules, and
  we have a natural bijection
  between irreducible $A$-modules
  and irreducible $B$-modules
  by $W_i \leftrightarrow U_i$.
\end{lemma}

\begin{proof}[Proof of Theorem \ref{thm:schur-weyl}]
  (1) Let $A$ be the image of
  $U(\gl_n)$ in $\End(V^{\otimes N})$
  and $B$ the image of
  $\C S_N$ in $\End(V^{\otimes N})$.
  By Lemma \ref{lem:double-centralizer},
  $A$ is a centralizer $Z_B$ of $B$,
  so $Z_B = S^N(\End(V))$. This proves
  part (1).

  (2) Since $\lambda$ has $\le n$
  parts, we have $L_\lambda \subseteq V^{\otimes N}$.

  (3) Since $\dim V \ge N$,
  the vectors $v_1, \dots, v_N$ are
  linearly independent. Then
  $\C S_N \to V^{\otimes N}$
  by $s \mapsto s(v_1 \otimes \cdots \otimes v_N)$
  is injective, so $B \cong \C S_N$.
\end{proof}
